\section{Discussion}
\label{sec:discussion}

The three experimental phases converge on a finding that is simultaneously specific---it concerns puzzle hunt authoring on a 45-point weighted rubric---and general: domain labels tell an AI what kind of output to produce, while artistic orientation labels tell it to ask what the person on the other end will feel before deciding on output format at all. We discuss the implications of this distinction and its structural supports.

\subsection{Domain Labels Versus Artistic Orientation}

The central finding of this study is a contrast between two cognitive stances activated by two single-word persona labels. ``Puzzle designer'' activates a domain template: conventions, formats, mechanism types, and quality criteria accumulated by practitioners of the field. ``Artist'' activates an experience-design orientation: the disposition to ask ``what does this person do, and how does it feel?'' before asking ``what does this artifact contain?''

This is not a puzzle-specific finding. A domain label tells the AI what kind of output to produce---the format, the genre, the conventions that practitioners of the domain would recognize as correct. An artistic label tells the AI to make experience-design the primary constraint, before format, genre, and convention are decided. For creative tasks, the second question is more fundamental, because format and genre conventions are means rather than ends. They exist to serve the experience of the person who will encounter the artifact. When they are activated as ends in themselves---which is what a domain label does---they constrain the generative space to the conventional solutions the domain has already produced.

The evidence for this interpretation is in the mechanism data. Every non-artist condition in the isolation study (T1 through T6, A0, and implicitly the puzzle designer control referenced in the introduction) remained within the letter-extraction mechanism family. The artist condition alone escaped it. The escape was not produced by domain knowledge that the artist persona encoded---the artist persona encodes no puzzle knowledge---but by the experience-first orientation that asked ``what does the solver do?'' before asking ``what letters spell the answer?'' The answer to the first question was: reconstruction of damaged material, which is structurally isomorphic to the experiential act of deciphering something ancient. That answer could not be reached by starting with the conventions of puzzle hunt design, because those conventions do not include damaged-inscription reconstruction. It required designing from the solver's phenomenological position outward, rather than from the mechanism family inward.

\subsection{A Hypothesis on Mechanism Family Escape}

The artist condition was the only persona to escape the letter-extraction mechanism family (acrostic, indexing) that dominated all other conditions. We propose the following hypothesis: experience-first orientation (T1) reframes the design problem before conventional mechanism templates activate. Under a domain label (``puzzle designer''), the model's problem representation is ``produce an AoE2 puzzle'' --- a representation that immediately activates the extraction family as the canonical puzzle hunt mechanism. Under an artistic orientation, the problem representation shifts to ``produce an experience of discovery from damaged material'' --- a representation for which the letter-extraction family is not the obvious instantiation. This is analogous to P\'erez y P\'erez's engagement/reflection phases in MEXICA~\cite{perezyperez2001mexica}: engagement (orientation-driven) determines which creative schemas are activated before reflection (evaluation) begins. T1 is the engagement operator; without it, generation proceeds directly to reflection over a default schema pool.

\subsection{The Three-Trait Injection and Its Generalization}

The minimum sufficient set (T1 + T4 + T5) has a practical form that can be stated in three sentences applicable to any creative brief across any domain: ``Think first about how the other person will experience this. Remove anything that is not necessary. Verify every claim before finishing.''

These three sentences do not encode puzzle design knowledge, AoE2 unit costs, or letter-extraction conventions. They encode a way of working. The first sentence reorients from content delivery to experience design. The second reorients from accumulation to reduction. The third reorients from confident assertion to verified claim. Together, they constitute the operational core of quality-oriented creative work in any domain where a human being will eventually encounter the output.

This generalization has theoretical grounding in the literature the traits operationalize. Experience-first thinking (T1) is the application of Norman's central design principle---that the user's mental model is the primary constraint, not the structural properties of the artifact---to the creative generation context~\cite{norman2013design}. Elegance drive (T4) is the application of Rams's ``as little design as possible'' to the prompt-level instruction~\cite{rams1976design}: the creator who removes anything that is not necessary is applying Rams's discipline not to a physical product but to a generated artifact. Rigor and verification (T5) is the implementation hygiene that keeps the other two honest. Without T5, experience-first orientation can produce beautiful-but-broken puzzles (the solver's experience is designed with care, but the clue that claims Watch Tower has a Blacksmith prerequisite is wrong, and the solver who knows this cannot solve the puzzle). Without T5, elegance drive can produce stripped-but-incorrect artifacts. The three traits are not independent quality contributions; they form a feedback loop. T1 sets the target (the solver's experience), T4 enforces discipline toward that target (remove what does not serve it), and T5 verifies that what remains is true (claims that are false serve no target).

We predict that a business analyst applying these three sentences to a deliverable will produce a more audience-oriented, less cluttered, and more accurate artifact than one who does not, for the same structural reasons the three-trait puzzle passes the C8 panel and the two-trait puzzle does not (predicted; cross-domain validation pending). The domain is irrelevant to this prediction; the cognitive stance is the operative variable.

\subsection{Why Elegance Drive Backfired in Isolation}

The T4 isolation result (26.0/45, below the null baseline of 24.0) is the study's most counterintuitive data point, and it deserves a direct explanation. The elegance drive trait, applied alone, produced the study's second-worst puzzle. The stripped-bracket presentation---five clues with the answer letters visible in brackets, no narrative, no misdirection---made the extraction mechanism so transparent that there was nothing for the solver to discover. The brackets presenting \texttt{[T]}, \texttt{[O]}, \texttt{[W]}, \texttt{[E]}, \texttt{[R]} in a vertical column give the solver the answer before a single clue is read. Stripping form without experience-first orientation strips the solver's journey along with the clutter.

The elegance paradox this reveals is precise: T4 alone produces transparency, not clarity. Clarity is a property of the relationship between the solver's understanding and the mechanism's structure---the mechanism is clear when the solver's journey to find it is unobstructed. Transparency is the elimination of that journey entirely. The difference between a clean puzzle and a transparent one is that the clean puzzle withholds until the solver earns the mechanism, while the transparent one removes the withholding along with the clutter.

This is why T4 needs T1 to know what to keep. Experience-first orientation identifies the things that serve the solver's journey and the things that obstruct it; elegance drive then removes the obstruction. Without T1, T4 cannot distinguish journey from clutter and removes both. The three minimum sufficient traits are not additive; T4 operates as a constraint on an experience-designed space that T1 defines, and T5 operates as a verification pass on the contents of that space. The traits form a dependency chain, not an independent feature list.

\subsection{Implications for AI Creative Pipeline Design}

The findings suggest three practical recommendations for practitioners deploying AI creative generation systems.

First, replace domain labels with artistic orientation labels for creative tasks. ``You are a puzzle designer,'' ``you are a marketing copywriter,'' ``you are a product specification author''---these labels activate domain conventions that constrain the generative space to what the domain has already produced. ``You are an artist who works in this domain'' or, more precisely, the three-sentence injection (T1, T4, T5)---activates the experience-design orientation that makes domain knowledge serve the audience rather than serve itself. The puzzle designer knows what an acrostic is; the artist asks whether an acrostic is what the solver needs.

Second, when domain expertise is genuinely required, inject it as content context rather than persona identity. The Phase~3 reconstruction experiments demonstrate this directly: the CALIBER puzzle (R3, T1+T4+T5) uses the game's actual values as the mechanism not because the artist persona encoded AoE2 knowledge but because the world data file was available as a resource. Domain expertise as content context (factual reference material, world files, knowledge bases) allows the model to draw on domain knowledge without adopting the domain's conventions as constraints on the design stance. The model with T1+T4+T5 and access to a world file asks ``how can this data serve the solver's experience?'' The model with a ``puzzle designer'' identity label asks ``how can I build a puzzle using this data?'' The first question produces CALIBER; the second produces The Indexed Garrison.

Third, the minimum three-sentence injection (T1+T4+T5) can be applied without domain-specific calibration. It does not reference puzzle design, AoE2, letter extraction, or any domain concept. It references only the cognitive stance toward quality that crosses the threshold in this study. We predict that practitioners who prepend these three sentences to any creative generation prompt---marketing brief, training material, interactive fiction specification, product design requirements---will observe a gap between domain-labeled and artistically-oriented output analogous to the 15.6-point finding documented here (predicted; cross-domain validation pending). The connection to Paper~\#3's finding~\cite{dellaLibera2025panel} is structural: that study demonstrates that authoring profile depth mirrors reviewing profile depth in its effect on output quality. The trait decomposition reported here explains why profile depth matters---richer profiles activate more of the six traits that constitute artistic quality orientation, and the traits that the richer profile activates more fully (T1, T4, T5) are precisely the ones whose absence causes the largest quality degradations.
